\section{Methodology}

Building on our observation that cross-benchmark correlation may reflect a shared latent factor, we design a factor-analytic framework with three components: (1) confound-controlled residualization, (2) latent factor extraction with permutation testing, and (3) mechanism validation through behavioral stability correlation.

\subsection{Confound Control via Residualization}

For each benchmark $b$, we regress scores on $\log(\text{parameters})$ and release date:
\begin{equation}
s_b = \beta_0 + \beta_1 \log_2(\text{params}) + \beta_2 \text{release\_date} + \epsilon_b
\end{equation}
The residuals $\epsilon_b$ represent benchmark performance unexplained by model scale or training recency.

\textbf{Data:} $N = 4{,}561$ models from the Open LLM Leaderboard with complete scores on 6 benchmarks (IFEval, BBH, MATH Lvl 5, GPQA, MUSR, MMLU-PRO).

\subsection{Latent Factor Extraction}

We apply Principal Component Analysis (PCA) to the $N \times 6$ residualized score matrix. To assess statistical significance, we construct a null distribution by shuffling model-benchmark pairings 1,000 times. The observed $\lambda_1$ is compared to the 95th percentile of this null distribution.

\subsection{Behavioral Stability Index (BSI)}

We construct a Behavioral Stability Index (BSI) measuring output consistency across semantically equivalent inputs, operationalized through paraphrase consistency on PAWS \citep{zhang2019paws} and QQP \citep{wang2018qqp} datasets.

\subsection{Instruction-Tuning Analysis}

We analyze 16 matched base/instruct model pairs to provide quasi-intervention evidence, testing whether instruction-tuning increases both BSI and PC1 scores.

\subsection{Prospective Validity}

We freeze PC1 weights and compute loadings on 6 holdout benchmarks (Truthfulness, Safety, Fairness, Robustness, Privacy, Ethics).
